% GENERATED by scripts/build-lecture-slides.py; do not edit. Portable profile: XeTeX, theme and fonts from the core.
\makeatletter\IfFileExists{core/beamer/beamerthememx.sty}{\def\mxcorepath{core/}\def\input@path{{core/beamer/}{core/}}}{\def\mxcorepath{../core/}\def\input@path{{../core/beamer/}{../core/}}}\makeatother
\documentclass[aspectratio=169,10pt,t]{beamer}
\usetheme{mx}
\usepackage{graphicx,booktabs}
\usepackage[french]{babel}
\setlength{\parskip}{.45em}
\hypersetup{pdfauthor={},pdfsubject={Lecture slides},pdfkeywords={},pdftitle={Équations et inéquations}}
\providecommand{\N}{\mathbb{N}}
\providecommand{\Z}{\mathbb{Z}}
\providecommand{\D}{\mathbb{D}}
\providecommand{\Q}{\mathbb{Q}}
\providecommand{\R}{\mathbb{R}}
\providecommand{\C}{\mathbb{C}}
\providecommand{\vect}[1]{\overrightarrow{#1}}
\providecommand{\vi}{\vec{\imath}}
\providecommand{\vj}{\vec{\jmath}}
\providecommand{\vk}{\vec{k}}
\providecommand{\norm}[1]{\left\lVert #1\right\rVert}
\providecommand{\vc}[2]{\begin{pmatrix}#1\\#2\end{pmatrix}}
\providecommand{\deter}[2]{\det\left(#1,#2\right)}
\providecommand{\abs}[1]{\left\lvert #1\right\rvert}
\providecommand{\dd}{\,\mathrm{d}}
\providecommand{\e}{\mathrm{e}}
\providecommand{\Card}{\operatorname{Card}}
\providecommand{\E}{\operatorname{E}}
\providecommand{\V}{\operatorname{V}}
\providecommand{\degre}{^{\circ}}
\providecommand{\un}[1]{\,\mathrm{#1}}
\providecommand{\sci}[2]{#1\times10^{#2}}
\providecommand{\pgcd}{\operatorname{PGCD}}
\providecommand{\ppcm}{\operatorname{PPCM}}
\title{Équations et inéquations}
\subtitle{Cours et exemples}
\date{}
\author{}
\begin{document}
\mxtitleframe
\begin{frame}{}
Un rectangle a pour côtés \(x+2\) et \(x+5\). Pour quelles valeurs de \(x\) son aire
vaut-elle \(28\)? Mettre le problème sous forme d'équation, chercher d'abord par
essais organisés, puis transformer l'expression pour faire apparaître un trinôme.
Comparer les stratégies et discuter les contraintes géométriques sur \(x\).
\end{frame}
\begin{frame}{Objectifs}
\begin{itemize}

  \item Résoudre les équations et inéquations du second degré.
  
\item Factoriser \(ax^2+bx+c\) grâce aux racines et à la forme canonique.
  
\item Étudier le signe d'un trinôme du second degré.
  
\item Résoudre des systèmes d'équations linéaires.

\end{itemize}
\end{frame}
\begin{frame}{Trinôme du second degré}
Un \textbf{trinôme du second degré} est une expression \(f(x) = ax^2+bx+c\)
avec \(a \ne 0\) et \(a,b,c \in \mathbb{R}\).
\end{frame}
\begin{frame}{Trinôme du second degré}
Pour \(ax^2+bx+c=0\) (\(a \ne 0\)), le \textbf{discriminant} est \(\Delta = b^2-4ac\).

\begin{itemize}

  \item \(\Delta < 0\) : aucune racine réelle, \(S = \emptyset\).
  
\item \(\Delta = 0\) : racine double \(x_0 = -\dfrac{b}{2a}\), 
        \(ax^2+bx+c = a(x-x_0)^2\).
  
\item \(\Delta > 0\) : deux racines
        \(x_{1,2} = \dfrac{-b \pm \sqrt{\Delta}}{2a}\), 
        \(ax^2+bx+c = a(x-x_1)(x-x_2)\).

\end{itemize}
\end{frame}
\begin{frame}{Trinôme du second degré}
Si \(x_1, x_2\) sont les deux racines :
\[x_1 + x_2 = -\dfrac{b}{a},\qquad x_1 x_2 = \dfrac{c}{a}.\]
\end{frame}
\begin{frame}{Signe du trinôme}
Si \(\Delta > 0\) (\(x_1 < x_2\)) et \(a > 0\) :

Si \(a > 0\) et \(\Delta \leq 0\) : \(ax^2+bx+c \geq 0\) sur tout \(\mathbb{R}\).
\end{frame}
\begin{frame}{Forme canonique}
Tout trinôme \(f(x)=ax^2+bx+c\) se met sous la forme :
\[f(x) = a\!\left(x + \dfrac{b}{2a}\right)^2 - \dfrac{\Delta}{4a},\]
où \(\Delta = b^2 - 4ac\). On pose \(\alpha = -\dfrac{b}{2a}\) (abscisse du sommet)
et \(\beta = -\dfrac{\Delta}{4a}\) (valeur au sommet).
\emph{Méthode pratique :} compléter le carré en extrayant le coefficient
de \(x^2\), puis ajouter et soustraire le bon terme.
\end{frame}
\begin{frame}{Forme canonique}
\(f(x) = 2x^2 - 8x + 3\).
\[f(x) = 2(x^2 - 4x) + 3 = 2\bigl[(x-2)^2 - 4\bigr] + 3
       = 2(x-2)^2 - 5.\]
Sommet : \((2, -5)\). Minimum de \(f\) : \(-5\), atteint en \(x=2\).
\end{frame}
\begin{frame}{Polynôme de degré 3}
Un \textbf{polynôme de degré 3} est une expression de la forme
\(P(x) = ax^3 + bx^2 + cx + d\) avec \(a \ne 0\).

On dit que \(r\) est une \textbf{racine} de \(P\) si \(P(r) = 0\).
\end{frame}
\begin{frame}{Polynôme de degré 3}
Si \(r\) est une racine de \(P(x) = ax^3+bx^2+cx+d\), alors il existe un
trinôme \(Q(x) = ax^2 + ex + f\) tel que :
\[P(x) = (x - r)\cdot Q(x).\]
On détermine \(Q\) par l'une des deux méthodes suivantes.
\end{frame}
\begin{frame}{Polynôme de degré 3}
On écrit \(P(x) = (x-r)(ax^2+ex+f)\) et on développe.
En identifiant les coefficients de chaque puissance de \(x\) avec ceux de \(P(x)\),
on trouve \(e\) et \(f\).

\emph{Exemple :} \(P(x) = x^3 - 6x^2 + 11x - 6\), racine \(r=1\).
\(P(x) = (x-1)(x^2+ex+f)\).
Développement : \(x^3 + ex^2 + fx - x^2 - ex - f
= x^3 + (e-1)x^2 + (f-e)x - f\).
Identification : \(e-1 = -6 \Rightarrow e=-5\) ;
\(-f = -6 \Rightarrow f=6\) ; vérif : \(f-e = 6-(-5)=11\).
Donc \(P(x) = (x-1)(x^2-5x+6)\).
\end{frame}
\begin{frame}{Polynôme de degré 3}
On divise \(P(x)\) par \((x-r)\) comme une division euclidienne :
\(P(x) = (x-r)\cdot Q(x) + 0\) (le reste est nul si \(r\) est racine).

\emph{Exemple :} \(P(x) = x^3 - 6x^2 + 11x - 6\) divisé par \((x-1)\).
\[
\begin{array}{r|l}
x^3-6x^2+11x-6 & x-1\\

-(x^3-x^2) & x^2-5x+6\\

\phantom{x^3{}}-5x^2+11x-6 & \\
\phantom{x^3{}}-(-5x^2+5x) & \\

\phantom{x^3-6x^2{}}6x-6 & \\
\phantom{x^3-6x^2{}}-(6x-6) & \\

\phantom{x^3-6x^2+11x{}}0 &
\end{array}
\]
Donc \(P(x) = (x-1)(x^2-5x+6) = (x-1)(x-2)(x-3)\).
\end{frame}
\begin{frame}{Polynôme de degré 3}
Si \(P(x) = a(x-r_1)(x-r_2)(x-r_3)\) avec \(r_1 < r_2 < r_3\) et \(a>0\) :
\end{frame}
\begin{frame}{Polynôme de degré 3}
\(P(x) \geq 0 \Leftrightarrow x \in [1,2]\cup[3,+\infty[\).
\end{frame}
\begin{frame}{Fractions rationnelles}
Une \textbf{fraction rationnelle} est une expression de la forme
\(F(x) = \dfrac{P(x)}{Q(x)}\) où \(P\) et \(Q\) sont des polynômes.
Son \textbf{ensemble de définition} est \(\mathbb{R} \setminus \{x \mid Q(x)=0\}\).
\end{frame}
\begin{frame}{Fractions rationnelles}
\begin{enumerate}
\setcounter{enumi}{0}
  \item Factoriser numérateur et dénominateur.
  
\item Simplifier les facteurs communs \emph{en notant les valeurs interdites}.

\end{enumerate}

\emph{Exemple :} \(F(x) = \dfrac{x^2-4}{x-2} = \dfrac{(x-2)(x+2)}{x-2} = x+2\)
pour \(x \ne 2\).
\end{frame}
\begin{frame}{Fractions rationnelles}
\begin{enumerate}
\setcounter{enumi}{0}
  \item Factoriser \(P(x)\) et \(Q(x)\).
  
\item Dresser le tableau de signe de chaque facteur.
  
\item Conclure : \(F(x)\) a le signe du produit num./dén.
        (attention : les zéros du dénominateur sont des \emph{valeurs interdites},
        pas des zéros de \(F\)).

\end{enumerate}
\end{frame}
\begin{frame}{Fractions rationnelles}
\(F(x) = \dfrac{(x+1)(x-3)}{x-1}\), définie sur \(\mathbb{R}\setminus\{1\}\).
\end{frame}
\begin{frame}{Systèmes d'équations}
Deux méthodes :

\begin{itemize}

  \item \textbf{Substitution} : exprimer une variable en fonction de l'autre
        et substituer.
  
\item \textbf{Combinaison linéaire} : multiplier les équations par des
        coefficients pour faire disparaître une variable.

\end{itemize}
\end{frame}
\begin{frame}{Systèmes d'équations}
\(\begin{cases}3x - y = 1 \\ x + 2y = 12\end{cases}\)

Multiplier la première par \(2\) : \(6x - 2y = 2\).
Ajouter la deuxième : \(7x = 14\), donc \(x=2\).
Retour : \(y = 3\times2-1 = 5\).
\end{frame}
\end{document}
